\chapter{Descriptive Statistics and Moments}
\label{cap:probabilidade}

This chapter presents \hayashi{}'s descriptive statistics tools as
expressions of probability theory. Each function maps directly to a
theoretical concept, making the code readable both for those coming from
mathematics and for those coming from Stata or R.

\section{Expected value}

The expected value $E(X)$ of a discrete random variable is the
weighted average of its values. For observed data, the natural estimator
is the sample mean:

\[
  \hat{E}(X) = \bar{X} = \frac{1}{n}\sum_{i=1}^{n} x_i
\]

In \hayashi{}, \lstinline{mean} accepts a list or a DataFrame column:

\begin{lstlisting}[caption={E(X) over a list and over a column}]
// over a list
let xs = [2.0, 4.0, 4.0, 4.0, 5.0, 5.0, 7.0, 9.0]
display mean(xs)    // 5.0

// over a DataFrame column
load "data.csv" as df
display mean(df, income)
\end{lstlisting}

\section{Variance and standard deviation}

The sample variance and sample standard deviation are:

\[
  \widehat{\mathrm{Var}}(X) = s^2 = \frac{1}{n-1}\sum_{i=1}^{n}(x_i - \bar{X})^2
  \qquad
  s = \sqrt{s^2}
\]

The divisor $n-1$ (Bessel's correction) yields an unbiased estimator of
$\sigma^2$. All \hayashi{} functions use this convention.

\begin{lstlisting}[caption={Variance and standard deviation}]
display variance(df, income)   // s²
display sd(df, income)         // s
\end{lstlisting}

\begin{nota}
  \lstinline{sd} and \lstinline{std} are aliases — both compute the sample
  standard deviation with divisor $n-1$.
\end{nota}

To obtain all moments at once, \lstinline{summarize} returns a dictionary:

\begin{lstlisting}
let stats = summarize(df, income)

display stats["mean"]      // mean
display stats["sd"]        // standard deviation
display stats["variance"]  // variance
display stats["min"]       // minimum
display stats["max"]       // maximum
display stats["N"]         // number of observations
\end{lstlisting}

\section{Conditional moments}

The conditional expected value $E(X \mid C)$ restricts the calculation to
a subset of the sample without creating an auxiliary DataFrame. The
\lstinline{if =} option is available in all scalar aggregation functions:

\begin{lstlisting}[caption={Conditional moments with if =}]
// E(income | sector == "public")
display mean(df, income, if = sector == "public")

// Var(income | sector == "public")
display variance(df, income, if = sector == "public")

// SD(income | sector == "public")
display sd(df, income, if = sector == "public")
\end{lstlisting}

The condition is evaluated row by row on the DataFrame — any Boolean
expression over columns is valid:

\begin{lstlisting}[caption={Compound conditions}]
display mean(df, wage,   if = educ >= 12 and sex == 1)
display sd(df,   income, if = state == "TX" or state == "CA")
\end{lstlisting}

\begin{dica}
  \lstinline{if =} does not modify the DataFrame. Use \lstinline{filter}
  to create a permanent subsample; use \lstinline{if =} for a one-off
  calculation.
\end{dica}

\section{Covariance}

The sample covariance between $X$ and $Y$ measures their joint variation:

\[
  \widehat{\mathrm{Cov}}(X, Y)
  = \frac{1}{n-1}\sum_{i=1}^{n}(x_i - \bar{X})(y_i - \bar{Y})
\]

\begin{lstlisting}[caption={Covariance between two columns}]
display cov(df, educ, wage)

// conditional: formal workers only
display cov(df, educ, wage, if = formal == 1)
\end{lstlisting}

\begin{nota}
  $\mathrm{Cov}(X, X) = \mathrm{Var}(X)$. To verify:
  \lstinline{cov(df, x, x)} returns the same value as
  \lstinline{variance(df, x)}.
\end{nota}

\section{Correlation}

\subsection{Scalar between two variables}

Pearson's correlation coefficient normalises the covariance by the product
of standard deviations:

\[
  \rho(X, Y) = \frac{\mathrm{Cov}(X,Y)}{\mathrm{SD}(X)\,\mathrm{SD}(Y)}
  \in [-1,\, 1]
\]

\lstinline{corr_pair} returns this scalar directly:

\begin{lstlisting}[caption={Scalar correlation}]
display corr_pair(df, educ, wage)

// conditional: private sector
display corr_pair(df, educ, wage, if = sector == "private")
\end{lstlisting}

\subsection{Correlation matrix}

To explore all relationships among multiple variables, use
\lstinline{correlate}:

\begin{lstlisting}[caption={Full correlation matrix}]
// all numeric variables
correlate(df)

// explicit subset
correlate(df, educ, exp, wage)

// with p-values (* p<0.1  ** p<0.05  *** p<0.01)
pwcorr(df, educ, exp, wage)
\end{lstlisting}

\begin{nota}
  \lstinline{correlate} prints the matrix but does not return a value —
  use \lstinline{corr_pair} when you need the scalar in a calculation.
\end{nota}

\section{Median and quantiles}

The median is the central value of the ordered distribution, less
sensitive to extreme values than the mean.

\begin{lstlisting}[caption={Median}]
display median(df, income)

// conditional
display median(df, income, if = region == "South")
\end{lstlisting}

Quantiles generalise the median to any point $p \in [0, 1]$ of the
distribution:

\begin{lstlisting}[caption={Quantiles}]
display quantile(df, income, 0.25)   // 1st quartile
display quantile(df, income, 0.50)   // median
display quantile(df, income, 0.75)   // 3rd quartile
display quantile(df, income, 0.90)   // 90th percentile
\end{lstlisting}

\lstinline{quantile} accepts \lstinline{if =}:

\begin{lstlisting}
// 90th percentile of income among college-educated workers
display quantile(df, income, 0.90, if = educ >= 16)
\end{lstlisting}

Both functions accept lists directly:

\begin{lstlisting}
let xs = [3.0, 1.0, 4.0, 1.0, 5.0, 9.0, 2.0, 6.0]
display median(xs)           // 3.5
display quantile(xs, 0.75)   // 5.25
\end{lstlisting}

\section{Detailed summary}

\lstinline{summarize} with \lstinline{detail=true} expands the return
dictionary to include skewness, kurtosis, and the full percentile grid:

\begin{lstlisting}[caption={summarize with detail}]
let s = summarize(df, income, detail=true)

display s["mean"]     // mean
display s["sd"]       // standard deviation
display s["p25"]      // 1st quartile
display s["p50"]      // median
display s["p75"]      // 3rd quartile
display s["p90"]      // 90th percentile
display s["p99"]      // 99th percentile
display s["skew"]     // skewness
display s["kurt"]     // kurtosis
\end{lstlisting}

\section{Quick reference}

\begin{center}
\begin{tabular}{lll}
\toprule
\textbf{Notation} & \textbf{Hayashi} & \textbf{Forms} \\
\midrule
$E(X)$
  & \lstinline{mean(df, x)}
  & list / df / \lstinline{if =} \\
$\mathrm{Var}(X)$
  & \lstinline{variance(df, x)}
  & list / df / \lstinline{if =} \\
$\mathrm{SD}(X)$
  & \lstinline{sd(df, x)}
  & list / df / \lstinline{if =} \\
$\mathrm{Cov}(X, Y)$
  & \lstinline{cov(df, x, y)}
  & df / \lstinline{if =} \\
$\rho(X, Y)$
  & \lstinline{corr_pair(df, x, y)}
  & df / \lstinline{if =} \\
$\mathrm{Med}(X)$
  & \lstinline{median(df, x)}
  & list / df / \lstinline{if =} \\
$Q_p(X)$
  & \lstinline{quantile(df, x, p)}
  & list / df / \lstinline{if =} \\
\midrule
\multicolumn{2}{l}{\lstinline{correlate(df, x, y, ...)}}
  & correlation matrix (display) \\
\multicolumn{2}{l}{\lstinline{summarize(df, x)}}
  & full dictionary of moments \\
\bottomrule
\end{tabular}
\end{center}

\section{Complete example}

The snippet below illustrates the joint use of these functions in a typical
exploratory analysis before estimating a model:

\begin{lstlisting}[caption={Exploratory analysis — income and education}]
load "survey.csv" as df

// overview
summarize(df, income)
summarize(df, educ)
correlate(df, income, educ, exp)

// moments by group
display mean(df, income, if = sex == 1)      // men
display mean(df, income, if = sex == 0)      // women

display sd(df, income, if = sex == 1)
display sd(df, income, if = sex == 0)

// gap between median and mean: sign of skewness
display mean(df, income)
display median(df, income)

// linear association
display cov(df, educ, income)
display corr_pair(df, educ, income)

// concentration at the top
display quantile(df, income, 0.90)
display quantile(df, income, 0.99)
\end{lstlisting}
